\section{Minecraft as a Construction Substrate}
\label{sec:minecraft-mechanics}

Because the mod's central claim is that building a circuit inside Minecraft is a
meaningfully different experience from drawing one in a schematic tool, it is worth
describing precisely what ``building'' means in Minecraft itself, and which of the game's
own mechanics the mod reuses, extends, or deliberately avoids.

\subsection{The voxel world and block placement}

Minecraft's world is a three-dimensional grid of one-metre cubes, or \emph{blocks}; almost
every action available to a player -- building a house, digging a mine, redirecting a
river -- reduces to placing and removing blocks from this grid. A block occupies exactly one
grid cell and, depending on its type, can carry a small amount of state directly in that
cell (its \emph{block state}): a small enumeration such as which of the six axis-aligned
directions it is \emph{facing}, whether a door is open, or how much liquid a cauldron holds.
When a player places a block that has a facing property, the game records the direction the
player was looking at the moment of placement as that block's initial facing; the block can
later be rotated by other in-game tools, but its facing is otherwise fixed until changed.
This mechanic -- an object whose orientation is set once, at placement time, and persists
until deliberately altered -- is exactly what every two-terminal component in this mod
reuses to determine which two of its six faces are its electrical leads
(Section~\ref{sec:architecture}): a resistor placed while facing north has its leads on its
north and south faces, precisely because that is already how the base game's own oriented
blocks (a piston, a dispenser, a chest with a facing hinge) behave, so no new placement
mechanic had to be introduced for players already familiar with the base game.

\subsection{Redstone: Minecraft's native circuit system}

The base game already ships an entire circuit-like subsystem, redstone, built around
\emph{redstone dust} (a wire-like block that carries a signal), \emph{redstone torches} and
other sources (which emit a signal), and a large family of blocks that read an incoming
signal to trigger some effect (opening a door, activating a piston, and so on). A redstone
signal is an integer strength from 0 to 15 that decreases by one for every block of dust it
travels through, and updates propagate through a discrete neighbor-notification event fired
once per game tick, rather than continuously; there is no built-in notion of voltage,
current, capacitance, resistance, or continuous time anywhere in the system. Despite this,
redstone is widely used, both informally and as a subject of academic
interest~\cite{nebel2016mining}, as an entry point into digital logic and even rudimentary
computer architecture -- players build AND/OR/NOT gates, adders, and working memory cells
entirely out of redstone components -- and it is common for players to describe and reason
about redstone circuits in explicitly electrical language even though the underlying
mechanic is discrete and boolean, not continuous~\cite{dezuanni2015redstone}. This is a
useful data point for this mod's design: it suggests that Minecraft's building metaphor
already primes players to think electrically, but that redstone itself cannot be the
vehicle for teaching \emph{analog} circuit behavior, since it has no continuous quantity to
teach with. Community-built ``analog redstone'' techniques exist that approximate a
continuous value by quantizing it into a comparator's 0--15 signal strength, but this only
ever yields a coarse, sixteen-level discretization, not a continuously varying, numerically
accurate quantity.

\subsection{Block entities and the tick loop}

A plain block's state is limited to the small, fixed set of properties declared for its
block type -- enough for a facing direction or an open/closed flag, but not for an arbitrary
amount of mutable data such as a capacitor's stored charge or 200 samples of recent voltage
history. Minecraft's modding interfaces expose a second, richer construct for exactly this
situation, a \emph{block entity}: an object associated with a specific block position that
can hold arbitrary data and, if the block type requests it, is given an opportunity to run
code once per game tick (20 times per second, the same rate the rest of the game simulates
at). Every component, wire node, ground reference, and configuration module in this mod is
implemented as a block entity for precisely this reason: it is the mechanism the platform
already provides for ``a block that needs to remember something and do something regularly,''
and reusing it means the mod's simulation cadence is naturally locked to the game's own tick
rate rather than requiring a separate timer mechanism.

\subsection{What the mod adds on top}

Given the above, the mod's own contribution is deliberately narrow: it introduces a second,
parallel notion of adjacency (``is this face electrically conductive toward that
neighbor'', Section~\ref{sec:architecture}) alongside the game's existing physical
adjacency, and it hands the resulting graph to a real solver every tick instead of hard-coding
any behavior into the blocks themselves. Redstone is reused, deliberately, only for the one
concept it is genuinely well suited to represent in this mod -- a boolean on/off switch
gating whether a voltage source is connected at all (Section~\ref{sec:pedagogy}) -- rather
than being pressed into service as an analog signal it was never designed to carry.
